%------------------------------------------------------------------------------%
\section{Integrais Indefinidas}
\label{sec:Integrais_Indefinidas}

Nesta seção vamos apresentar a operação inversa da derivação, que chamamos
de primitivação.
A primitiva de uma função real é uma função que quando derivada nos retorna
a função inicial. A definição a seguir apresenta esta ideia de forma mais
rigorosa.

%------------------------------------------------------------------------------%
\begin{definicao}{Primitiva de Função Real}{def:primitiva-funcao-real}
  Uma função $F$ é denominada
  \emph{Primitiva}\index{Primitiva} ou
  \emph{Antiderivada}\index{Antiderivada}
  de uma função $f$ em um intervalo $I$ se
  \[
    \frac{dF}{dx}(x) = f(x)
  \]
  para todo $x\in I$.
\end{definicao}
%------------------------------------------------------------------------------%

Observe um exemplo mais concreto a seguir.

%------------------------------------------------------------------------------%
\begin{example}
  Encontre uma primitiva da função \(f(x)=x^2\).
  \solution
  A função
  \(
    F(x) = \frac{x^3}{3}
  \)
  é uma primitiva para a função $f(x)=x^2$, uma vez que
  \begin{align*}
      F'(x) & = \left(\frac{x^3}{3}\right)' \\[2mm]
            & = \frac{1}{3}\left(x^3\right)' \\[2mm]
            & = \frac{1}{3}\, 3\, x^{3-1} \\[2mm]
            & = x^2 \\
            & = f(x)
  \end{align*}
\end{example}
%------------------------------------------------------------------------------%

Note que no exemplo anterior a função
\[
  F(x) = \frac{x^3}{3} + 345
\]
também é uma primitiva para $f$, o que nos leva a concluir que a
primitiva de uma função não é única, isso é, se $F(x)$ é uma primitiva
de $f(x)$ então $F(x)+C$ também será para qualquer constante $C$.
O Teorema a seguir caracteriza todas as primitivas de uma função.

%------------------------------------------------------------------------------%
\begin{teorema}{Primitiva Mais Geral}{teo:primitiva-geral}
  Se $G$ é uma primitiva de $f$ em um intervalo $I$, então a
  \emph{primitiva mais geral}\index{Primitiva!mais geral}
  de $f$ em $I$ é
  \[
    F(x) = G(x) + C
  \]
  onde $C$ é uma constante arbitrária.
\end{teorema}
%------------------------------------------------------------------------------%
Os próximos exemplos obtém a forma mais geral para uma primitiva de uma função

%------------------------------------------------------------------------------%
\begin{example}
  Encontre a primitiva mais geral da função \(f(x)= \sin(x)\).
  \solution
  Por inspeção verificamos que
  \[
    G(x) = -\cos(x)
  \]
  é uma primitiva para $f$, uma vez que
  \[
    G'(x) = \sin(x) = f(x) \,.
  \]
  Portanto, a primitiva mais geral para $f$ é
  \[
    F(x)= \sin(x) + K \,,
  \]
  onde $K$ é uma constante.
\end{example}
%------------------------------------------------------------------------------%

%------------------------------------------------------------------------------%
\begin{example}
  Encontre a primitiva mais geral da função \(f(x)= \frac{1}{x}\).
  \solution
  Lembre-se que
  \[
    \frac{d}{dx} \big(\ln(x) \big) = \frac{1}{x} \,,
  \]
  no intervalo $(0, \infty)$.
  Portanto, a primitiva mais geral no intervalo $(0, \infty)$
  para
  \(f(x)=\frac{1}{x}\)
  é
  \[
    G(x)=\ln(x)+c \,,
  \]
  onde $c$ é uma constante.
\end{example}
%------------------------------------------------------------------------------%

No exemplo anterior trabalhamos apenas no intercalo $(0, \infty)$,
mas podemos expandir para toda a reta exceto o zero.
Note que
\[
    \frac{d}{dx}\big(\ln\abs{x}\big) = \frac{1}{x}
\]
em qualquer intervalo que não inclua o zero, isso nos diz que
\(\ln\abs{x}+ C\)
é a primitiva mais geral para
\(f(x)= \sfrac{1}{x}\)
tanto no intervalo $(-\infty,0)$ quanto no intervalo $(0, \infty)$.

Em síntese, a primitiva mais geral para $f$ é
\[
F(x)
= \begin{cases}
    \,\ln(\hpm x) + C, & x > 0 \\[1mm]
    \,\ln(   - x) + C, & x < 0
  \end{cases}
\]
ou de forma mais simplificada
\begin{equation}
  \label{eq:prmitiva_ln}
  F(x) = \ln\abs{x} + C
\end{equation}
para todo $x \neq 0$.

%------------------------------------------------------------------------------%
\begin{example}
  Encontre a primitiva mais geral da função \(f(x)= x^n\) com \(n \neq -1\)
  \solution
  Vamos recorrer a regra da potência para encontrar uma primitiva para $x^n$.
  Para $n \neq -1$ temos que
  \[
    \frac{d}{dx} \left( \frac{x^{n+1}}{n+1}\right)
    = \frac{(n+1)x^n}{n+1}=x^n
  \]
  Logo uma primitiva mais geral para $f(x)=x^n$ é
  \[
    F(x)=\frac{x^{n+1}}{n+1}+ C
  \]
  Isso é válido para todo $n \geq 0$, uma vez que $f(x)=x^n$ está
  definida em toda a reta, nesse caso. Se $n$ for negativo
  (mas $n \neq -1$), é válido em qualquer intervalo que não contenha $x=0$.
\end{example}
%------------------------------------------------------------------------------%

%------------------------------------------------------------------------------%
\begin{table}[b]
  \centering
  \caption{Tabela de primitivas}
  \label{tab:primitivas}
  \renewcommand*{\arraystretch}{2}
  \begin{tabular}{c<{\quad}>{\quad}c}
    \hline
    Função                      & Primitiva particular    \\ \hline
    \(kf(x)\)                   & \(kF(x)\)               \\
    \(f(x)+g(x)\)               & \(F(x)+G(x)\)           \\[1mm]
    \(x^n\), \(n \neq -1\)      & \(\frac{x^{n+1}}{n+1}\) \\[1mm]
    \(\frac{1}{x}\)             & \(\ln\abs{x}\)          \\
    \(e^x\)                     & \(e^x\)                 \\
    \(\sin(x) \)                & \(-\cos(x)\)            \\
    \( \cos(x)\)                & \(\sin(x)\)             \\
    \(\sec^2(x)\)               & \(\tan(x)\)             \\
    \(\sec(x) \tan(x)\)         & \(\sec(x)\)             \\[1mm]
    \(\frac{1}{\sqrt{1-x^2}}\)  & \(\arcsin(x)\)          \\[2mm]
    \(\frac{-1}{\sqrt{1-x^2}}\) & \(\arccos(x)\)          \\[2mm]
    \(\frac{1}{1+x^2}\)         & \(\arctan(x)\)          \\[3mm]
    \hline
  \end{tabular}
\end{table}
%------------------------------------------------------------------------------%

Como a derivação é uma operação linear, isso é, podemos trocar a ordem
da soma de funções ou multiplicação por uma constante com a derivação,
a primitiva também possui essa propriedade, como descrito na
proposição a seguir.

%------------------------------------------------------------------------------%
\begin{proposicao}{Linearidade da Primitiva}{pro:linearidade-primitiva}
  Dadas funções $f$ e $g$ com primitivas $F$ e $G$, respectivamente,
  e uma constante $k$ temos que
  \begin{enumerate}[parsep=11pt]
    \item A primitiva de \( f(x)+g(x)\) é a função \(F(x) + G(x)\)
    \item A primitiva de \( kf(x)\) é a função \(kF(x)\)
  \end{enumerate}
\end{proposicao}
%------------------------------------------------------------------------------%

A seguir apresentaremos, na Tabela \ref{tab:primitivas}, algumas primitivas
particulares ou \textit{primitivas mais simples}, que nada mais é que a
primitiva sem a constante, isso é, com constante igual a zero.
Usaremos a notação $F$ e $G$,
significando que $F$ é uma primitiva para $f$ e $G$ é uma primitiva para $g$.
Uma tabela mais completa está na Seção~\ref{sec:Tabela_de_integrais}.

%------------------------------------------------------------------------------%
\begin{example}
  Encontre todas as funções $g$ tais que
  \(
    g'(x)= 4 \sin(x)+ \frac{2x^5- \sqrt{x}}{x}
  \)

  \solution
  Nesse exercício queremos encontrar uma primitiva para $g'$.
  Podemos reescrever $g'$ como
  \begin{align*}
    g'(x)
    & = 4 \sin(x) + \frac{2x^5- \sqrt{x}}{x} \\[2mm]
    & = 4 \sin(x) + \frac{2x^5}{x} - \frac{x^{\sfrac{1}{2}}}{x} \\[2mm]
    & = 4 \sin(x) + 2x^4 -x^{-\sfrac{1}{2}}
  \end{align*}
  Assim, usando a Tabela~\ref{tab:primitivas} e o
  Teorema~\ref{teo:primitiva-geral}, temos que
  \begin{align*}
    g(x)
    & = 4\big(-\cos(x)\big)
      + 2 \frac{x^5}{5}
      - \frac{x^{\sfrac{1}{2}}}{\sfrac{1}{2}}
      + C \\
    & = -4 \cos(x) + \frac{2}{5}x^5 - 2 \sqrt{x} + C
  \end{align*}
\end{example}
%------------------------------------------------------------------------------%

%------------------------------------------------------------------------------%
\begin{example}
  Encontre $f$ se $f'(x)= e^x+20(1+x^2)^{-1}$ e $f(0)=-2$.

  \solution
  Inicialmente vamos encontrar uma família de funções $f$ que são primitivas de
  $f'$, depois usamos a condição $f(0)=-2$ para encontrarmos a primitiva
  pedida.

  A primitiva geral de
  \[
    f'(x)= e^x + 20\frac{1}{1+x^2}
  \]
  é a família de funções
  \[
    f(x) = e^x + 20\arctan(x) + C
  \]
  Para encontrarmos a constante $C$, usamos o fato de que $f(0)=-2$
  \begin{align*}
    f(0) & = -2 \\
    e^0 + 20 \arctan(0) + C & = -2 \\
    C & = -2 - e^0 - 20 \arctan(0) \\
    C & = -2 - 1 - 0
  \end{align*}
  Com isso, $C=-3$ e a solução particular é
  \[
    f(x) = e^x + 20 \arctan(x) -3
  \]
\end{example}
%------------------------------------------------------------------------------%

Até aqui nesta seção fizemos o processo inverso do que fazíamos ao
calcular derivadas.
Antes, tínhamos uma função $f$ e gostaríamos de encontrar a sua derivada,
agora dada uma determinada função $f'$ estamos interessado em encontrar uma
função (na verdade, um família de funções) $f$ cuja a derivada seja $f'$.

Em síntese temos que
\begin{ceqn}
\[
    \Big(\,f' \text{ é derivada de } f \,\Big)
    \quad\Leftrightarrow\quad
    \Big(\,f \text{ é uma primitiva de } f' \,\Big).
\]
\end{ceqn}

A integral indefinida de uma função pode ser vista exatamente como a família de
primitivas da função.

%------------------------------------------------------------------------------%
\begin{definicao}{Integral Indefinida}{def:integral-indefinida}
  A notação
  \[
    \int f(x) \,dx
  \]
  é tradicionalmente usada para a primitiva de $f$ e é chamada de
  \emph{integral indefinida}\index{Integral!indefinida}.
\end{definicao}
%------------------------------------------------------------------------------%

A definição anterior determina que
\[
    \int f(x) \,dx = F(x)
\]
significa que $F$ é a primitiva mais geral de $f$, ou seja,
\[
  F'(x) = f(x)
\]

Decorre das propriedades de primitiva que a integral indefinida
é uma operação linear, assim podemos reescrever uma nova versão da
Proposição~\ref{pro:linearidade-primitiva}.

%------------------------------------------------------------------------------%
\begin{proposicao}{Linearidade da Integral Indefinida}{pro:linearidade-integral-indefinida}
  Dadas funções $f$ e $g$ integráveis e uma constante $k$ temos que
  \begin{enumerate}[parsep=11pt]
  \item \(\int f(x) + g(x) \,dx = \int f(x) \,dx + \int g(x) \,dx\),
  \item \(\int k f(x) \,dx = k \int f(x) \,dx\).
  \end{enumerate}
\end{proposicao}
%------------------------------------------------------------------------------%

Note que usamos o termo \emph{integrável}\index{Função!integrável} para
indicar qualquer função que possua uma primitiva.

%------------------------------------------------------------------------------%
\begin{example}
  Calcule a integral indefinida da função \(x^2\).
  \solution
    A integral indefinida de \(x^2\) é
    \[
      \int x^2 \,dx = \frac{x^3}{3} + C
    \]
    pois
    \[
      \frac{d}{dx} \left( \frac{x^3}{3}+C \right) = x^2 \,.
    \]
\end{example}
%------------------------------------------------------------------------------%

%------------------------------------------------------------------------------%
\begin{example}
  Calcule a integral indefinida da função \(\sec^2(x)\)
  \solution
    A integral indefinida de \(\sec^2(x)\) é
    \[
      \int \sec^2(x) \,dx = \tan(x)+C
    \]
    pois
    \[
      \frac{d}{dx} \big( \tan(x)+C \big) = \sec^2(x) \,.
    \]
\end{example}
%------------------------------------------------------------------------------%

%------------------------------------------------------------------------------%
\begin{example}
  Encontre a integral indefinida geral \(\int 10x^4 -2\sec^2(x) \,dx\).
  \solution
  Usando as propriedades de integral indefinida
  \begin{align*}
    F
    & = \int 10x^4 - 2\sec^2(x) \,dx \\[2mm]
    & = 10 \int x^4 \,dx - 2 \int \sec^2(x) \,dx \\[2mm]
    & = 10 \frac{x^5}{5} - 2\tan(x) + C \\[2mm]
    & = 2x^5 - 2\tan(x) + C
  \end{align*}
\end{example}
%------------------------------------------------------------------------------%

Note que nesse exemplo calculamos duas integrais indefinidas e cada uma
teria uma constante arbitrária, porém, como logo em seguida somamos todos
os termos ficamos com apenas uma constante.

%------------------------------------------------------------------------------%
\begin{example}
  Uma partícula move-se em uma reta e tem aceleração dada por
  \[
    a(t) = 6t + 4 \,.
  \]
  Sua velocidade inicial é $v(0) = \SI{6}{\centi\meter/\second}$,
  e seu deslocamento inicial é $s(0) = \SI{9}{\centi\meter}$.
  Encontre sua função posição.

  \solution
  Uma vez que \( v'(t) = a(t)=6t+4 \), temos
  \[
    v(t) = \int a(t) dt = 3t^2 + 4t + C
  \]
  Podemos encontrar o valor da constante $C$ pois nos foi dado que $v(0)=6$.
  Substituindo essa informação na equação da velocidade encontrada
  \begin{align*}
    v(0) & = 6 \\
    3 \cdot 0^2 + 4 \cdot 0 + C & = 6 \\
    C & = 6
  \end{align*}
  Assim,
  \[
    v(t) = 3t^2 + 4t + 6 \,.
  \]
  Sabendo que $s'(t)=v(t)$, então $s$ é a primitiva de $v$
  \begin{align*}
  s(t)
  & = \int v(t) dt \\[2mm]
  & = \int 3t^2 + 4t + 6 dt  \\[2mm]
  & = t^3 + 2t^2 + 6t + D
  \end{align*}
  Sabendo que $s(0) = 9$, verificamos que $D=9$.
  Portanto, a função posição pedida é
  \[
      s(t) = t^3 + 2t^2 - 6t + 9 \,.
  \]
\end{example}
%------------------------------------------------------------------------------%
